% 本文件由 LaTeX 排版 Agent 根据有效 translation.json 自动生成。
% author=Councils; work_id=3808; title=君士坦丁堡第一届大公会议（公元三八一年）

\chapter{君士坦丁堡第一届大公会议（公元三八一年）}

% structure: p0001 source=h1 target=omitted reason=page-title-already-rendered-by-parent

% structure: p0002 source=h2 target=section reason=relative-heading-rank; base=h2

\section{信经}

我们信唯一的天主，全能的圣父，天地万物，无论有形无形，都是祂所创造的。我们信唯一的主耶稣基督，天主的独生子，在万世之前，由圣父所生，出自光明的光明，出自真天主的真天主，受生而非受造，与圣父同一性体；万物藉着祂而受造。祂为了我们人类，并为我们的救恩，从天降下，由圣神和童贞玛利亚降生成人，成为人；祂在般雀比拉多执政时为我们被钉在十字架上，受难而被埋葬，第三天按照圣经复活了，并升了天，坐在圣父的右边。祂还要光荣地降来，审判生者死者；祂的王国万世无终。（\Emph{I}）\par

我们信圣神，主及赋予生命者，由圣父所发；祂与圣父、圣子同受钦崇，同享光荣；祂曾藉先知发言。我们信唯一、圣、（\Emph{II}）至公、宗徒的教会。我们承认一个圣洗圣事以赦免罪过；我们期待死人的复活及来世的生命。阿们。\par

% structure: p0005 source=h2 target=section reason=relative-heading-rank; base=h2

\section{同一神圣主教会议致最虔诚的皇帝狄奥多西大帝的信函，附其所颁之教会法规}

致最虔诚的皇帝狄奥多西，自不同省份聚集在君士坦丁堡的主教神圣会议。\par

我们以对天主的感谢开始致陛下虔敬的信，祂为教会的共同和平及对真信仰的支持建立了陛下的帝国。在向天主献上应尽的感谢之后，我们依本分向陛下陈述神圣会议中已完成之事。因此，当我们按照陛下之书函聚集在君士坦丁堡时，首先我们彼此更新了心灵的合一，然后我们颁布了一些简洁的定义，确认了尼西亚神长们的信仰，并绝罚了与之相反的兴起异端。此外，我们还为更好地管理教会订定了一些教会法规，全部附于本信之后。因此，我们恳请陛下批准主教会议的谕令，以便您既以召书尊崇了教会，也请为决议的结论盖上印章。愿主在和平与正义中巩固陛下的帝国，并将其延绵万代；并愿祂将天上王国的享福也加于陛下的地上权柄。愿天主藉\OriginalTerm{圣人们的祈祷}{\GreekText{εὐχαῖς} \GreekText{τῶν} \GreekText{ἁγίων}}眷顾世界，使您作为最虔诚、蒙天主所爱的皇帝，在一切美事上强大而卓越。\par

% structure: p0008 source=h2 target=section reason=relative-heading-rank; base=h2

\section{一百五十位神长在君士坦丁堡集会期间、在光耀者弗拉维乌斯·尤切里乌斯与弗拉维乌斯·埃瓦格里乌斯执政时期、七月七日前之第七日所颁之教会法规}

因天主的恩宠，应最虔诚的皇帝狄奥多西之召，聚集在君士坦丁堡的不同省份的主教们议决如下：\par

% structure: p0010 source=h3 target=subsection reason=relative-heading-rank; base=h2

\subsection{第一道法令}

在尼西亚（比提尼亚）聚集的三百一十八位神长所确立的信仰不得作废，应常保坚固。凡异端均应受绝罚，尤其尤诺米派或（\Emph{阿诺米派}、亚流派或）尤多克西派，以及半亚流派或敌圣神派，撒伯流派，马尔塞流派，弗提诺派，阿波利纳流派。\par

% structure: p0012 source=h3 target=subsection reason=relative-heading-rank; base=h2

\subsection{第二道法令}

主教不得越界至其教区之外的教会，以免扰乱教会；但亚历山大里亚主教应依教会法规单独管理埃及事务；东方主教应单独管理东方事务，同时保存安提约基雅教会在尼西亚法规中所提及的特权；亚细亚教区主教应仅管理亚细亚事务；本都主教仅管理本都事务；色雷斯主教仅管理色雷斯事务。主教不得为圣职授予或其他教会职务而越界，除非受邀请。若上述关于教区的法规得到遵守，则显然各省的主教会议将依尼西亚所定管理该省事务。但在异教民族中的天主教会应按自教父时代以来所通行的惯例管理。\par

% structure: p0014 source=h3 target=subsection reason=relative-heading-rank; base=h2

\subsection{第三道法令}

关于犬儒者马克西穆及因此人在君士坦丁堡发生的混乱，兹决定：马克西穆从未是、现在也不是主教；凡由他所祝圣的人在圣职中均无任何品级；因为凡为他或由他所行之事，均宣布为无效。\par

% structure: p0016 source=h3 target=subsection reason=relative-heading-rank; base=h2

\subsection{第四道法令}

（可能于次年即三八二年在君士坦丁堡举行的一次会议中通过。见：关于教会法规数目的导论。）\par

关于西方［主教］的卷册，我们在安提约基雅也接纳那些宣认圣父、圣子及圣神同一神性的人。\par

% structure: p0019 source=h3 target=subsection reason=relative-heading-rank; base=h2

\subsection{第五道法令}

凡从异端归向正统信仰，并归入得救行列者，我们按以下方法与惯例接纳：亚略派、马其顿派、撒巴提乌斯派、诺瓦提安派（他们自称卡塔里派或阿里斯托利派）、夸托得辛派（或称泰特拉迪特派）、以及阿波林派，我们接纳他们，条件是他们书面弃绝〔其谬误〕，并诅咒每个不符合天主的神圣、公教、宗徒教会的异端。此后，他们首先被用圣油在额、眼、鼻、口、耳上盖印或傅油；当我们盖印时，我们说：\Quote{圣神恩赐的印记。}但欧诺米派（他们仅以一次浸礼受洗）、孟他努派（此处称为弗里吉亚派）、撒伯流派（他们教导圣父与圣子同一位格，并做其他各种有害之事）、以及所有其他异端——因为此处有许多这类人，尤其来自加拉提亚地区者——所有这些，当他们愿意归向正统信仰时，我们视同外教人接纳。第一日我们使他们成为基督徒；第二日成为慕道者；第三日我们向他们面部和耳朵三次嘘气以驱魔；然后我们教导他们，并责成他们在教会内度过一段时间，聆听圣经；然后给他们施洗。\par

\SourceRecord{Translated by Henry Percival.；Nicene and Post-Nicene Fathers, Second Series；Vol. 14.；Edited by Philip Schaff and Henry Wace.；Buffalo, NY: Christian Literature Publishing Co.,；1900.；<http://www.newadvent.org/fathers/3808.htm>.}
